%% kapa: exam=1 round=37 year=2026 period=1 form=A subject=real_estate subject_ko=부동산학원론 seq=39 number=119
%% source: Q-Net 감정평가사 기출문제 (cst00309) artlSeq=5258924 file=raw/1008/5258924/2026년 제37회 감정평가사 1차 1교시.pdf
%% conversion: kordoc 4.24.51 -> tools/gen1.py (auto); review=proofread-sonnet; flags=none
\begin{question}{119}
거래사례비교법에 관한 내용으로 옳지 않은 것은?
\begin{choices}
  \item 거래사례비교법이란 대상물건과 가치형성요인이 같거나 비슷한 물건의 거래사례와 비교하여 대상물건의 현황에 맞게 사정보정, 시점수정, 가치형성요인 비교 등의 과정을 거쳐 대상물건의 가액을 산정하는 감정평가방법을 말한다.
  \item 시점수정은 대상물건의 가격 변동률로 함이 원칙이다.
  \item 감정평가법인등이 토지를 감정평가하는 경우 적정한 실거래가가 있는 경우에는 이를 기준으로 할 수 있다.
  \item 인근지역에 적정한 거래사례가 없는 경우 동일수급권 안의 유사지역에 존재하는 거래사례를 적용할 수 있다.
  \item 거래사례와 대상물건 간에 종별ㆍ유형별 특성에 따라 지역요인이나 개별요인 등 가치형성요인에 차이가 있는 경우에는 이를 각각 비교하여 대상물건의 가치를 개별화ㆍ구체화하여야 한다.
\end{choices}
\end{question}
\answer{2}
